% MA 214 Discussion 4 -- covers Lectures 6 and 7.
%
% Student handout. Page 1 is the concept review; the question set runs from page 2.
% Compile from inside this folder:  pdflatex Discussion4.tex

\documentclass[11pt]{article}
\usepackage[margin=1in]{geometry}
\usepackage{parskip}
\usepackage{fancyhdr}
\usepackage{array}
\usepackage{booktabs}
\usepackage{enumitem}
\usepackage{amsmath}
\usepackage{tabularx}
\usepackage{listings}

\pagestyle{fancy}
\fancyhf{}
\cfoot{\thepage}
\rhead{Discussion 4}
\lhead{MA 214: Applied Statistics}
\renewcommand{\headrulewidth}{.4mm}

\newcommand{\blankline}[1]{\underline{\hspace{#1}}}
\newcolumntype{L}{>{\raggedright\arraybackslash}X}

\lstset{
  basicstyle=\ttfamily\small,
  columns=fullflexible,
  frame=single,
  framerule=0.4pt,
  breaklines=true,
  showstringspaces=false,
  aboveskip=8pt,
  belowskip=8pt
}

% Section header: bold title over a hairline rule.
\newcommand{\parttitle}[1]{%
  \par\medskip\noindent
  {\large\bfseries #1}\par
  \vspace{-.5em}\noindent\rule{\linewidth}{0.4pt}\par\smallskip}

% Writing space.
\newcommand{\answerspace}[1]{\vspace{#1}}

% Flush-left question list: the label runs in at the margin and the body wraps
% back to the margin, so nothing is pushed far to the right.
\newlist{questions}{enumerate}{1}
\setlist[questions]{label=\textbf{Question \arabic*.}, wide=0pt, itemsep=1.4em, topsep=.8em}

\newlist{parts}{enumerate}{1}
\setlist[parts]{label=\textbf{(\Alph*)}, leftmargin=1.6em, labelsep=.45em,
  itemsep=.7em, topsep=.5em}

\begin{document}

\noindent\textbf{Name:}\blankline{2.3in}\hfill\textbf{BUID:}\blankline{1.6in}

\begin{center}
  {\Large\bfseries Discussion 4: Odds Ratios and Choosing a Model}
\end{center}

\noindent\textbf{Goals:} \textbf{(1)} read a logistic coefficient as an odds ratio,
\textbf{(2)} move between log odds, odds, and probability, \textbf{(3)} use AIC to choose among
candidate models, \textbf{(4)} recognize what multicollinearity does to a coefficient table,
\textbf{(5)} explain why accuracy can be worthless on unbalanced data.

\parttitle{Part 1: Concept Review}

{\small

\noindent Everything in Part 2 can be answered from this page. Keep it in front of you.

\begin{enumerate}[label=\textbf{\arabic*.}, wide=0pt, itemsep=.5em, topsep=.5em]

\item \textbf{Three scales for the same thing.} With $p = P(Y = 1)$,
\[
\text{odds} = \frac{p}{1-p},
\qquad
\text{log odds} \;=\; \eta \;=\; \log\!\frac{p}{1-p},
\qquad
p = \frac{e^{\eta}}{1+e^{\eta}} .
\]
Logistic regression is linear in the \emph{log odds}:
$\eta = \beta_0 + \beta_1 x_1 + \cdots + \beta_p x_p$. It is not linear in $p$, which is why the
same coefficient moves the probability by different amounts at different starting points.

\item \textbf{Coefficients are odds ratios.} Exponentiating a coefficient gives the multiplicative
effect on the \emph{odds}, holding the other predictors fixed:
$\text{OR} = e^{\beta}$ for a one-unit increase, and $e^{c\beta}$ for a $c$-unit increase. An
$\text{OR} > 1$ raises the odds, $\text{OR} < 1$ lowers them, and
$100(e^{\beta}-1)\%$ is the percent change in the odds. A confidence interval for the odds ratio
is found by exponentiating the endpoints of the interval for $\beta$, not the other way around.

\item \textbf{Fitting.} There is no least squares here. The model is fit by maximum likelihood on
the Bernoulli likelihood $\prod_i p_i^{y_i}(1-p_i)^{1-y_i}$, solved numerically.

\item \textbf{From a probability to a decision.} Choose a threshold, predict $1$ above it, and
cross-tabulate against the truth:
\[
\text{accuracy} = \frac{TP+TN}{n},
\qquad
\text{sensitivity} = \frac{TP}{TP+FN},
\qquad
\text{specificity} = \frac{TN}{TN+FP}.
\]
On \textbf{unbalanced} data, always predicting the majority class already scores an accuracy equal
to the majority share, so accuracy alone cannot show a model is useful. Lowering the threshold
raises sensitivity and lowers specificity; the right threshold depends on which error costs more.

\item \textbf{Choosing a model with AIC.} There is no $R^2$ for logistic regression, so use
\[
\text{AIC} = -2\,\ell(\hat\theta) + 2k,
\]
where $k$ counts the estimated parameters. \textbf{Lower is better.} The first term rewards fit,
the second charges for complexity. Differences below about 2 are not meaningful.

\item \textbf{Multicollinearity.} When two predictors carry nearly the same information, the data
cannot separate their individual contributions. Standard errors inflate, so $t$ or $z$ statistics
shrink and $p$-values grow, and each predictor can look useless while the pair jointly matters.
What it does \emph{not} damage is the fitted values or predictions. Interpretation breaks;
prediction survives.

\item \textbf{Four traps.} Reading $e^{\beta}$ as a ratio of probabilities, or as a probability.
Reporting accuracy without the base rate beside it. Comparing AIC across models fit to
\emph{different} data sets or different numbers of rows. Concluding a predictor does not matter
from a large $p$-value when a collinear twin is sitting in the same model.

\end{enumerate}

}

\newpage

\parttitle{Part 2: Question Set}

\noindent\textbf{Scenario.} For each of $n = 200$ bike-share stations you record whether it
\textbf{ran empty} at least once during the day ($1$ = yes), along with \texttt{docks} (docks in
tens, ranging 4 to 12) and \texttt{type} (Campus, Downtown, Residential). In the data,
\textbf{32 of the 200 stations ran empty}. Fitting \texttt{empty} $\sim$ \texttt{docks} $+$
\texttt{type} by logistic regression gives

\begin{center}
\renewcommand{\arraystretch}{1.2}
\begin{tabular}{lcccc}
\toprule
 & \textbf{Estimate} & \textbf{Std. Error} & \textbf{$z$} & \textbf{$p$-value} \\
\midrule
(Intercept) & $0.1137$ & 0.7074 & $0.16$ & 0.872 \\
\texttt{docks} & $-0.3152$ & 0.0927 & $-3.40$ & 0.0007 \\
\texttt{type}Downtown & $1.3892$ & 0.5085 & $2.73$ & 0.0063 \\
\texttt{type}Residential & $-0.3101$ & 0.6299 & $-0.49$ & 0.623 \\
\bottomrule
\end{tabular}
\qquad AIC $= 157.23$
\end{center}

\noindent You may use: $e^{-0.3152} = 0.730$, \; $e^{1.3892} = 4.012$, \;
$e^{-0.3101} = 0.733$, \; $e^{-1.7775} = 0.169$, \; $e^{-0.3883} = 0.678$, \;
$e^{-1.6492} = 0.192$.

\begin{questions}

%%%%%%%%%%%%%%%%%%%%%%%%%%%%%%%%%%%%%%
\item \textbf{Coefficients as odds ratios.}

\begin{parts}
\item Convert the \texttt{docks} coefficient to an odds ratio, and interpret it in context.
Remember that one unit of \texttt{docks} is ten docks.
\answerspace{4.5em}

\item Express the same effect as a percent change in the odds.
\answerspace{3em}

\item Interpret the \texttt{type}Downtown odds ratio. Your sentence must say what the comparison
group is, and must include the conditional phrase.
\answerspace{4em}

\item A classmate writes: ``Downtown stations are 4.01 times as likely to run empty.'' That is not
what $4.012$ says. Correct the sentence.
\answerspace{3.5em}

\item The 95\% confidence interval for the \texttt{docks} coefficient is $[-0.497,\; -0.134]$.
Give the 95\% interval for its odds ratio, and say why the interval does not contain 1.
\answerspace{4em}
\end{parts}

%%%%%%%%%%%%%%%%%%%%%%%%%%%%%%%%%%%%%%
\item \textbf{From log odds to a probability.} Use the fitted model above.

\begin{parts}
\item Fill in the table for a station with \texttt{docks} $= 6$.

\begin{center}
\renewcommand{\arraystretch}{1.7}
\begin{tabular}{lccc}
\toprule
\textbf{Type} & \textbf{Log odds $\eta$} & \textbf{Odds} & \textbf{Probability $p$} \\
\midrule
Campus & \blankline{1.2in} & \blankline{.9in} & \blankline{.9in} \\
Downtown & \blankline{1.2in} & \blankline{.9in} & \blankline{.9in} \\
\bottomrule
\end{tabular}
\end{center}

\item A Downtown station with \texttt{docks} $= 10$ has $\eta = -1.6492$. Compute its probability.
\answerspace{3em}

\item Going from 6 to 10 tens of docks at a Downtown station, the log odds changed by a fixed
$4 \times (-0.3152)$. Did the \emph{probability} drop by a fixed amount too? Compare your answers
in (A) and (B) and explain what this shows about the model.
\answerspace{4.5em}

\item \texttt{type}Residential has $p$-value $0.623$. Does that mean Residential stations run
empty at the same rate as Campus stations? Answer carefully.
\answerspace{4em}
\end{parts}

%%%%%%%%%%%%%%%%%%%%%%%%%%%%%%%%%%%%%%
\item \textbf{Choosing among four models.} All four are fit to the same 200 stations. The variable
\texttt{capacity} is the station's total capacity, which is essentially \texttt{docks} measured in
single docks rather than tens.

\begin{center}
\renewcommand{\arraystretch}{1.35}
\begin{tabular}{llcc}
\toprule
\textbf{Model} & \textbf{Predictors} & \textbf{$k$} & \textbf{AIC} \\
\midrule
1 & \texttt{docks} & 2 & 167.35 \\
2 & \texttt{docks} $+$ \texttt{type} & 4 & 157.23 \\
3 & \texttt{docks} $+$ \texttt{type} $+$ \texttt{capacity} & 5 & 158.91 \\
4 & \texttt{capacity} $+$ \texttt{type} & 4 & 157.52 \\
\bottomrule
\end{tabular}
\end{center}

\begin{parts}
\item Which model does AIC select? By how much does it beat the runner-up, and is that difference
meaningful?
\answerspace{4em}

\item Model 1 to Model 2 costs two extra parameters and AIC still improves by more than 10. What
does that tell you about \texttt{type}?
\answerspace{3.5em}

\item Models 2 and 4 have almost identical AIC. Explain why, using what \texttt{capacity} is. Is
this a case of two competing explanations, or something else?
\answerspace{4.5em}

\item Model 3 contains everything Models 2 and 4 contain, yet its AIC is the worst of the three.
How is that possible?
\answerspace{4em}
\end{parts}

%%%%%%%%%%%%%%%%%%%%%%%%%%%%%%%%%%%%%%
\item \textbf{What collinearity did to the table.} Here is Model 3, the one with both
\texttt{docks} and \texttt{capacity}. In this data set their correlation is $0.998$.

\begin{center}
\renewcommand{\arraystretch}{1.2}
\begin{tabular}{lcccc}
\toprule
 & \textbf{Estimate} & \textbf{Std. Error} & \textbf{$z$} & \textbf{$p$-value} \\
\midrule
(Intercept) & $0.1028$ & 0.7073 & $0.15$ & 0.884 \\
\texttt{docks} & $-1.1484$ & 1.4617 & $-0.79$ & 0.432 \\
\texttt{type}Downtown & $1.4029$ & 0.5099 & $2.75$ & 0.0059 \\
\texttt{type}Residential & $-0.2926$ & 0.6311 & $-0.46$ & 0.643 \\
\texttt{capacity} & $0.0833$ & 0.1456 & $0.57$ & 0.568 \\
\bottomrule
\end{tabular}
\end{center}

\begin{parts}
\item Compare the \texttt{docks} row here with the \texttt{docks} row in Model 2. By what factor
did its standard error change, and what happened to its $p$-value?
\answerspace{3.5em}

\item In Model 3, \emph{both} \texttt{docks} and \texttt{capacity} have large $p$-values. Has the
number of docks stopped mattering? Explain what the data cannot do here.
\answerspace{4.5em}

\item The \texttt{capacity} coefficient is \emph{positive} while the \texttt{docks} coefficient is
strongly negative, even though more capacity and more docks are the same thing. Why should you not
try to interpret either number?
\answerspace{4.5em}

\item What would you do about it, and what part of the analysis would be unaffected by the problem?
\answerspace{4em}
\end{parts}

%%%%%%%%%%%%%%%%%%%%%%%%%%%%%%%%%%%%%%
\item \textbf{Accuracy, and find the bugs.} A classmate evaluates Model 2 at a threshold of $0.5$
and gets this table.

\begin{center}
\renewcommand{\arraystretch}{1.3}
\begin{tabular}{lcc}
\toprule
 & \textbf{Actually ran empty} & \textbf{Did not run empty} \\
\midrule
Predicted ``ran empty'' & 3 & 3 \\
Predicted ``did not'' & 29 & 165 \\
\bottomrule
\end{tabular}
\end{center}

\begin{parts}
\item Compute the accuracy, the sensitivity, and the specificity.
\answerspace{4.5em}

\item Compute the accuracy of the rule ``predict that no station ever runs empty.'' Compare it with
your answer in (A) and say what that comparison establishes.
\answerspace{4em}

\item The city wants to find stations to send more bikes to. Which of the three numbers in (A)
matters most to them, and which way should the threshold move?
\answerspace{4em}

\item Their code is below. They intended to fit the logistic model from Question 2 and report its
accuracy. It runs with no error and prints a plausible number, and two separate lines need fixing.
Find both and give the correction.

\begin{lstlisting}[language=R]
fit <- glm(empty ~ docks + type, data = st)

phat <- predict(fit, newdata = st)

mean((phat > 0.5) == st$empty)
\end{lstlisting}
\answerspace{5.5em}
\end{parts}

\end{questions}

\end{document}
